\chapter{Teething}

\section*{Definition}
Teething is the process of primary (baby) teeth emerging through the gums. It typically begins around 6 months of age and continues until about 2-3 years. While teething can cause mild discomfort, it does not cause significant illness.

\section*{When to See a Doctor}
\begin{itemize}
\item Seek medical advice for a temperature of 38 C (100.4 F) or higher, persistent diarrhea or vomiting, poor fluid intake, lethargy, breathing difficulty, or symptoms that are severe or not explained by a tooth emerging.
\item Ear pulling, marked irritability, or disturbed sleep should not automatically be attributed to teething; assess for infection when appropriate.
\end{itemize}

\section*{How It Develops}
As a tooth moves through the jawbone and gum tissue, it causes localized inflammation and pressure. This triggers the release of inflammatory mediators that can cause pain, swelling, and tenderness of the gums. The cutting of the gum tissue is a mechanical process that naturally resolves once the tooth erupts.

\section*{Causes}
\begin{itemize}
\item Normal eruption of primary teeth
\item Typically begins at 6 months of age
\item Lower central incisors usually erupt first
\item Molars (around 12-18 months) may cause more discomfort
\end{itemize}

\section*{Related Symptoms}
\begin{itemize}
\item Gum discomfort, drooling, irritability, and sleep disturbance may occur, but fever and diarrhea are not explained by teething
\item Diarrhea (not caused by teething)
\item Rash on the face (from drooling)
\item Ear infection (can be confused with teething pain)
\end{itemize}

\section*{Medical Evaluation}
\begin{itemize}
\item Clinical diagnosis based on age and symptoms
\item Visual examination of the gums and emerging tooth
\end{itemize}

\section*{Treatment Options}
\begin{itemize}
\item Pain relief: acetaminophen or ibuprofen (age-appropriate dosing)
\item Teething rings and other non-pharmacologic measures
\item Do not use benzocaine or topical lidocaine for teething pain in infants and young children; benzocaine should not be used for teething in children under 2.
\item Avoid amber necklaces (no evidence of benefit and pose strangulation risk)
\end{itemize}

\section*{Self-Care and Home Management}
\begin{itemize}
\item Gently massage the gums with a clean finger
\item Provide a chilled (not frozen) teething ring
\item Offer cold, soft foods (yogurt, pureed fruit) for older infants
\item Use a clean, wet washcloth that has been chilled
\item Wipe drool frequently to prevent facial rash
\item Ask a healthcare professional about weight-based acetaminophen dosing; ibuprofen is generally avoided in infants younger than 6 months unless specifically advised.
\end{itemize}

\section*{References}
\begin{thebibliography}{99}
\bibitem{teething:ref1} Massignan C, Perazzo MF, Cardoso AMR, et al. ``Signs and symptoms of primary tooth eruption: A meta-analysis.'' \textit{Pediatrics}, 2016;137(3):e20153501.
\bibitem{teething:ref2} Macknin ML, Piedmonte M, Jacobs J, et al. ``Symptoms associated with infant teething: A prospective study.'' \textit{Pediatrics}, 2000;105(4):747--752.
\bibitem{teething:ref3} Wake M, Hesketh K, Lucas J. ``Teething and tooth eruption in infants: A cohort study.'' \textit{Pediatrics}, 2000;106(6):1374--1379.
\bibitem{teething:ref4} Tsang AKL. ``Teething in infants: Current practice and recommendations.'' \textit{Journal of the Canadian Dental Association}, 2010;76:a75.
\bibitem{teething:ref5} Markman L. ``Teething: Facts and fiction.'' \textit{Pediatric Review}, 2009;30(8):e59--e64.
\bibitem{teething:ref6} Ashley MP. ``Is it teething?'' \textit{Archives of Disease in Childhood}, 2001;85(2):149--150.
\end{thebibliography}
